\SourceLeaf{430}%
Fläche selbst gelegt wird. Für den Krümmungshalbmesser haben wir also die einfache Formel
\[
\frac1r=Z(\xi\xi T+2\xi\eta U+\eta\eta V).
\]

\begin{center}11.\end{center}
\noindent\hspace*{1.2em}Da sich durch diese Normale unendlich viele Ebenen legen lassen, so können hiernach unendlich viele verschiedene Werthe des Krümmungshalbmessers stattfinden. In dieser Beziehung sind $T,U,V,Z$ als constant, $\xi,\eta,\zeta$ als veränderlich zu betrachten. Um die letztern von Einer veränderlichen Grösse abhängig zu machen, nehme man in dem grössten Kreise, dessen Pol $L$ ist, zwei feste um $90^{0}$ von einander entfernte Punkte an, die $M,M'$ heissen mögen. Es seien ihre Coordinaten gegen den Mittelpunkt der Kugel $\alpha,\beta,\gamma$; $\alpha',\beta',\gamma'$. Man hat dann
\[
\cos\lambda(1)=\cos\lambda M.\cos M(1)+\cos\lambda M'.\cos M'(1)+\cos\lambda L.\cos L(1).
\]
Setzt man nun
\[
\lambda M=\varphi,
\]
so ist
\[
\cos\lambda M'=\sin\varphi
\]
und die Formel wird
\[
\xi=\alpha\cos\varphi+\alpha'\sin\varphi
\]
und eben so
\[
\begin{aligned}
\eta&=\beta\cos\varphi+\beta'\sin\varphi\\
\zeta&=\gamma\cos\varphi+\gamma'\sin\varphi.
\end{aligned}
\]

Also, wenn man
\[
\begin{aligned}
A&=(\alpha\alpha T+2\alpha\beta U+\beta\beta V)Z\\
B&=(\alpha\alpha'T+(\alpha\beta'+\alpha\beta')U+\beta\beta'V)Z\\
C&=(\alpha'\alpha'T+2\alpha'\beta'U+\beta'\beta'V)Z
\end{aligned}
\]
setzt:
